\section{Task 1B: Stride Pattern Prefetcher (SPP)}
\label{sec:task1b}
\subsection{Data structures and operation}
The supplied \texttt{spp.l1d\_pref} has three direct-mapped structures: 24 IP trackers, a 1,024-entry signature table, and a 512-entry prefetch tracker. The IP key is again masked to eight bits. Each IP tracker stores a 64-bit IP field, last cache-line address, last stride, 32-bit LRU rank, four signed 64-bit stride-history values, a history count, a confidence integer, and a valid flag. When a new IP is seen, an invalid entry is preferred; otherwise an LRU entry is reused and its history is reset.

For each nonzero stride, the implementation updates a short stride-history buffer. After four previous strides are available, it forms a 64-bit signature using the recurrence $s_{i+1}=31s_i+h_i$ over the four signed history values, then uses \texttt{signature \% 1024} as the index. A signature-table entry stores the full signature tag, the next observed stride, a confidence integer, and a valid flag. Matching observations increment confidence up to 7; a changed predicted stride replaces the stored stride and resets confidence to 1. A pattern-table prediction is accepted only when the full signature matches and its confidence is at least 2.

The predictor prioritizes a constant stride whenever its per-IP confidence is at least 1. That confidence is incremented when the current stride equals the previous stride and reset to 0 on a mismatch. If uniform confidence is 0 and four-history state is available, it attempts a signature-table prediction. Thus the source implements a compact conditional fallback, not a separate multi-state classifier with hysteresis. On each demand access it generates up to the current degree of next-stride requests; after each predicted request, the local history is shifted and the next prediction is selected again. A negative candidate or a cross-page candidate terminates the loop. Prefetches fill L1 when MSHR occupancy is below half capacity and L2 otherwise.

\subsection{Accuracy-based degree control}
Accepted prefetch requests increment the global \texttt{issued} counter. A prefetch-table entry, indexed by cache-line address modulo 512, records the address and a state: 1 for an accepted request not yet marked as filled, 2 after a matching prefetch fill, and 0 when cleared. A later demand access that hits a state-2 tracked line increments \texttt{useful}. On each interval where at least 256 additional requests have been accepted, the policy computes the usefulness ratio from the deltas since the last update and sets degree to 1 below 30\%, 2 from 30\% to below 60\%, and 3 at 60\% or above. Degree begins at 3. This is a batch update after issued requests, not a continuously updated sliding window.

\begin{table}[H]\centering\small
\caption{Degree selection encoded in the supplied SPP implementation.}
\begin{tabular}{p{0.45\linewidth}p{0.10\linewidth}p{0.37\linewidth}}
\toprule Measured accuracy since previous update & Degree & Effect \\
\midrule
Below 30\% & 1 & Conservative; limits traffic after low observed usefulness \\
30\% to below 60\% & 2 & Intermediate aggressiveness \\
60\% or above & 3 & Maximum degree in this implementation \\
\bottomrule
\end{tabular}\end{table}

\subsection{Logical storage accounting}
The accounting below counts declared field information bits, treating a boolean as one logical bit and a C++ \texttt{int} as 32 bits. The IP field is counted as 64 bits because the source stores it in a \texttt{uint64\_t}, even though the value is masked to eight bits before use.
\begin{table}[H]\centering\scriptsize
\caption{Bit-level logical storage for all SPP state. C++ object padding is reported separately below.}
\resizebox{\linewidth}{!}{\begin{tabular}{llrrr}
\toprule Structure & Field & Bits per entry & Entries & Total bits \\
\midrule
IP tracker & IP tag storage & 64 & 24 & 1,536 \\
IP tracker & Last cache-line address & 64 & 24 & 1,536 \\
IP tracker & Last stride & 64 & 24 & 1,536 \\
IP tracker & LRU rank & 32 & 24 & 768 \\
IP tracker & History[4] & 256 & 24 & 6,144 \\
IP tracker & History count & 32 & 24 & 768 \\
IP tracker & Confidence & 32 & 24 & 768 \\
IP tracker & Valid flag (logical bit) & 1 & 24 & 24 \\
\midrule
Signature entry & Signature tag & 64 & 1,024 & 65,536 \\
Signature entry & Predicted stride & 64 & 1,024 & 65,536 \\
Signature entry & Confidence & 32 & 1,024 & 32,768 \\
Signature entry & Valid flag (logical bit) & 1 & 1,024 & 1,024 \\
\midrule
Prefetch tracker & Cache-line address & 64 & 512 & 32,768 \\
Prefetch tracker & State integer & 32 & 512 & 16,384 \\
\midrule
Global state & issued, useful, last\_issued, last\_useful & 256 & 1 & 256 \\
Global state & Degree integer & 32 & 1 & 32 \\
\midrule
\multicolumn{4}{r}{\textbf{Total logical field bits}} & \textbf{227,384} \\
\multicolumn{4}{r}{\textbf{Logical bytes (bits/8)}} & \textbf{28,423} \\
\multicolumn{4}{r}{\textbf{Logical KiB (bytes/1024)}} & \textbf{27.76} \\
\bottomrule\end{tabular}}\end{table}

On a conventional 64-bit ABI, typical structure alignment would make the arrays approximately $24\times80 + 1024\times24 + 512\times16=34,688$ bytes. Adding the four 64-bit counters and the 32-bit degree gives about 34,724 bytes (33.91 KiB), excluding executable/runtime overhead. The logical-field total and ABI-size estimate answer different questions: the former counts information field widths, while the latter includes expected padding. Exact physical bytes should be confirmed with \texttt{sizeof} on the compiler used for the final submission.

\subsection{Measured results}
\begin{table}[H]\centering\scriptsize
\caption{Task 1B measurements. Coverage and accuracy are not applicable to the no-prefetch baseline.}
\resizebox{\linewidth}{!}{\begin{tabular}{llrrrrr}
\toprule Trace & Configuration & IPC & Speedup & L1D Load MPKI & Accuracy (\%) & Coverage (\%) \\
\midrule
Tr1.gz & No prefetch + LRU & 0.232345 & 1.000000 & 10.6780 & N/A & N/A \\
Tr1.gz & IP-stride & 0.233452 & 1.004764 & 10.1190 & 52.71 & 5.24 \\
Tr1.gz & SPP & 0.232577 & 1.000999 & 10.5465 & 41.33 & 1.23 \\
Tr2.gz & No prefetch + LRU & 0.251158 & 1.000000 & 13.5379 & N/A & N/A \\
Tr2.gz & IP-stride & 0.251260 & 1.000406 & 13.4378 & 38.18 & 0.74 \\
Tr2.gz & SPP & 0.251094 & 0.999745 & 13.5014 & 33.54 & 0.27 \\
Tr3.gz & No prefetch + LRU & 0.382065 & 1.000000 & 3.4034 & N/A & N/A \\
Tr3.gz & IP-stride & 0.382862 & 1.002086 & 3.3182 & 15.99 & 2.51 \\
Tr3.gz & SPP & 0.382139 & 1.000194 & 3.3938 & 8.51 & 0.28 \\
\bottomrule\end{tabular}}
\end{table}

\plotpair{plots/task1b/ipc.png}{IPC}{plots/task1b/speedup.png}{Speedup}{Task 1B IPC and speedup across traces}{fig:t1b-ipc-speedup}
\plotpair{plots/task1b/l1d_load_mpki.png}{L1D load MPKI}{plots/task1b/prefetch_accuracy_pct.png}{Prefetch accuracy}{Task 1B miss intensity and prefetch accuracy}{fig:t1b-mpki-accuracy}
\plotpair{plots/task1b/l1d_coverage_pct.png}{L1D coverage}{plots/task1b/ipc_change_vs_baseline_pct.png}{IPC change vs. baseline}{Task 1B coverage and relative IPC change}{fig:t1b-coverage-change}
\plotone{plots/task1b/l1d_mpki_reduction_vs_baseline_pct.png}{L1D load-MPKI reduction relative to no-prefetch + LRU. Positive values mean fewer L1D load misses.}{fig:t1b-mpki-reduction}

\subsection{Interpretation}
On Tr1, IP-stride improves IPC by 0.476\% over no-prefetch + LRU, while SPP improves it by 0.100\%. On Tr2, IP-stride improves by 0.041\%, while SPP is 0.025\% below baseline. On Tr3, IP-stride improves by 0.209\% and SPP by 0.019\%. Thus IP-stride has the highest IPC of the compared configurations on all three traces in this data set.

IP-stride also has higher measured accuracy and coverage than SPP on every trace. The L1D load-MPKI reductions are 5.24\%, 0.74\%, and 2.51\% for IP-stride on Tr1/Tr2/Tr3, compared with 1.23\%, 0.27\%, and 0.28\% for SPP. These observations do not prove that multi-stride prediction is generally worse; they show that the current compact signature/fallback implementation does not outperform the provided IP-stride baseline on these particular traces. The signature table's direct mapping, small IP tracker pool, masked IP key, confidence fallback, and modulo-indexed prefetch tracker are plausible limits, but the measurements alone do not isolate which mechanism explains the gap.
